\documentclass[12pt,letterpaper,oneside]{book}

\usepackage[utf8]{inputenc}
\usepackage[T1]{fontenc}
\usepackage[spanish,es-tabla]{babel}
\usepackage{csquotes}
\usepackage{newtxtext,newtxmath}
\usepackage{setspace}
\usepackage{geometry}
\usepackage{graphicx}
\usepackage{booktabs}
\usepackage{longtable}
\usepackage{array}
\usepackage{float}
\usepackage{caption}
\usepackage{subcaption}
\usepackage{xcolor}
\usepackage{hyperref}
\usepackage{cleveref}
\usepackage{titlesec}
\usepackage{tocloft}
\usepackage{calc}

% Lo que necesita el LaTeX que genera pandoc desde los capítulos en Markdown
\providecommand{\tightlist}{\setlength{\itemsep}{0pt}\setlength{\parskip}{0pt}}
\providecommand{\pandocbounded}[1]{#1}
% Símbolos Unicode que pdflatex no trae declarados
\DeclareUnicodeCharacter{2212}{\ensuremath{-}}
\DeclareUnicodeCharacter{2264}{\ensuremath{\leq}}
\DeclareUnicodeCharacter{2265}{\ensuremath{\geq}}
\DeclareUnicodeCharacter{2248}{\ensuremath{\approx}}
\DeclareUnicodeCharacter{2192}{\ensuremath{\rightarrow}}
\DeclareUnicodeCharacter{03A3}{\ensuremath{\Sigma}}
\DeclareUnicodeCharacter{221A}{\ensuremath{\surd}}

\geometry{
  letterpaper,
  left=3cm,
  right=2.5cm,
  top=3cm,
  bottom=3cm
}

\onehalfspacing
\setlength{\parindent}{1.25cm}
\setlength{\parskip}{0pt}

\hypersetup{
  colorlinks=true,
  linkcolor=black,
  citecolor=black,
  urlcolor=blue,
  pdftitle={Categorizacion de los impactos de variabilidad climatica global en el rendimiento de cereales},
  pdfauthor={Paola Andrea Ruiz Franco}
}

\usepackage[
  backend=biber,
  style=apa,
  natbib=true,
  sorting=nyt,
  language=spanish
]{biblatex}
\DeclareLanguageMapping{spanish}{spanish-apa}
\addbibresource{bib/referencias.bib}

\title{Categorización de los impactos de variabilidad climática global en el rendimiento de cereales y la generación de insumos para el desarrollo de seguros paramétricos orientados a fortalecer la seguridad alimentaria}
\author{Paola Andrea Ruiz Franco}
\date{2026}

\begin{document}

\frontmatter
\begin{titlepage}
  \centering
  \vspace*{1cm}

  {\Large \textbf{Categorización de los impactos de variabilidad climática global en el rendimiento de cereales y la generación de insumos para el desarrollo de seguros paramétricos orientados a fortalecer la seguridad alimentaria}\par}

  \vspace{2.5cm}

  {\large Paola Andrea Ruiz Franco\par}

  \vspace{2cm}

  {\large Tesis presentada como requisito parcial para optar al título de:\par}
  \vspace{0.3cm}
  {\large \textbf{Magíster en Ingeniería - Recursos Hidráulicos}\par}

  \vspace{2cm}

  {\large Director:\par}
  \vspace{0.3cm}
  {\large Carlos David Hoyos Ortiz\par}

  \vspace{0.6cm}

  {\large Codirector:\par}
  \vspace{0.3cm}
  {\large Santiago Medina Hurtado\par}

  \vfill

  {\large Universidad Nacional de Colombia\par}
  {\large Facultad de Minas\par}
  {\large Medellín, Colombia\par}
  {\large 2026\par}
\end{titlepage}

\chapter*{Dedicatoria}
\addcontentsline{toc}{chapter}{Dedicatoria}

% TODO: completar dedicatoria personal.

\vspace*{4cm}

\begin{flushright}
\textit{A quienes han acompañado este proceso académico y personal.}
\end{flushright}

% % Agradecimientos: pendiente de redactar.
% Para incluirlos, quitar el comentario de % Agradecimientos: pendiente de redactar.
% Para incluirlos, quitar el comentario de % Agradecimientos: pendiente de redactar.
% Para incluirlos, quitar el comentario de \input{frontmatter/agradecimientos} en main.tex.

\chapter*{Agradecimientos}
\addcontentsline{toc}{chapter}{Agradecimientos}
 en main.tex.

\chapter*{Agradecimientos}
\addcontentsline{toc}{chapter}{Agradecimientos}
 en main.tex.

\chapter*{Agradecimientos}
\addcontentsline{toc}{chapter}{Agradecimientos}
  % pendiente de redactar
\chapter*{Resumen}
\addcontentsline{toc}{chapter}{Resumen}

La seguridad alimentaria global depende de la estabilidad productiva de cultivos básicos como maíz, arroz, trigo y soya. Sin embargo, la variabilidad climática asociada a modos de gran escala, como El Niño-Oscilación del Sur (ENSO) y la Oscilación Madden-Julian (MJO), puede alterar la disponibilidad hídrica, la temperatura y la ocurrencia de eventos extremos durante fases críticas del desarrollo agrícola. Esta tesis propone caracterizar la relación entre la variabilidad climática global y el rendimiento de cereales, con el fin de identificar regímenes agroclimáticos y umbrales climáticos aplicables al diseño de seguros paramétricos agrícolas.

El documento se estructura en dos artículos científicos. El primero analiza la sensibilidad espacio-temporal del rendimiento de maíz, arroz, trigo y soya frente a ENSO y MJO, integrando anomalías climáticas, rendimientos agrícolas y calendarios fenológicos globales. El segundo utiliza los insumos derivados del primer artículo para proponer una metodología de identificación de umbrales climáticos y diseño preliminar de disparadores paramétricos orientados a la gestión del riesgo agrícola y al fortalecimiento de la seguridad alimentaria.

% TODO: completar con metodología definitiva, fuentes de datos, periodo de análisis y principales resultados propios.

\textbf{Palabras clave:} seguridad alimentaria; ENSO; MJO; cereales; rendimiento agrícola; seguros paramétricos; riesgo climático; fenología agrícola.

\chapter*{Abstract}
\addcontentsline{toc}{chapter}{Abstract}

Global food security depends on the production stability of staple crops such as maize, rice, wheat, and soybean. However, climate variability associated with large-scale modes such as the El Niño-Southern Oscillation (ENSO) and the Madden-Julian Oscillation (MJO) can modify water availability, temperature, and the occurrence of extreme events during critical crop development stages. This thesis aims to characterize the relationship between global climate variability and cereal yields in order to identify agroclimatic regimes and climate thresholds that may support the design of agricultural parametric insurance.

The thesis is organized around two scientific papers. The first paper analyzes the spatial and temporal sensitivity of maize, rice, wheat, and soybean yields to ENSO and MJO variability, integrating climate anomalies, crop yield data, and global crop calendars. The second paper uses the outputs from the first paper to propose a methodology for identifying climate thresholds and preliminary parametric triggers for agricultural risk management and food security applications.

% TODO: complete with final methodology, datasets, study period, and original results.

\textbf{Keywords:} food security; ENSO; MJO; cereals; crop yield; parametric insurance; climate risk; crop phenology.


\tableofcontents
\listoffigures
\listoftables

\mainmatter
% Los capítulos se escriben en chapters/*.md; convertir_capitulos.sh genera chapters/tex/*.tex antes de compilar.
\input{chapters/tex/01_introduccion}
\input{chapters/tex/02_marco_teorico}
\input{chapters/tex/03_datos}
\input{chapters/tex/04_metodologia}
\input{chapters/tex/05_resultados}
\input{chapters/tex/06_discusion}
\input{chapters/tex/07_seguros_parametricos}
\input{chapters/tex/08_conclusiones}

\backmatter
\chapter*{Bibliografía completa}
\addcontentsline{toc}{chapter}{Bibliografía completa}
\printbibliography[heading=none]

\end{document}
